\section{Learnable paired-feature polyconvex energies}
\label{sec:learned}
The non-intrusive route developed here replaces the repeated online solution of the microscopic problem by a learned effective stored energy. Its central constitutive relation is
\begin{equation}
 \ssv(\ee)=\nabla_{\ee}W(\ee),
 \label{eq:energy_stress_relation}
\end{equation}
where $\ssv$ is the stress work-conjugate to the engineering-strain vector $\ee$. Thus all stress components derive from a common potential, and the Hessian of $W$ supplies a major-symmetric material tangent. Fitting stress components independently does not, in general, enforce the cross-derivative identities required by such a potential~\cite{asad2022}. This observation motivates the energy-based formulation adopted throughout this section.

One could stop at this point and construct a flexible hyperelastic network without restricting the curvature of its learned potential. We use precisely this option as the \emph{Unconstrained energy} reference. It takes objective strain measures as inputs and obtains stress through Eq.~\eqref{eq:energy_stress_relation}. Beyond potential structure and objectivity, the only mechanical conditions imposed by construction are the reference conditions $W(\I)=0$ and $\partial_{\ee}W(\bm0)=\bm0$. The first fixes the arbitrary additive energy constant; the second is enforced by an affine correction that cancels the stress at the undeformed state, following the consistency construction of As'ad et al.~\cite{asad2022}.\footnote{No convexity or growth condition is imposed on this reference.}

The Unconstrained energy model shows what can be represented before stronger finite-strain restrictions are introduced. Potential structure and reference normalization alone do not control deformation-gradient curvature or prevent finite energy as $J\to0^+$. We therefore pursue polyconvexity by constructing the energy as a convex function of independent deformation minors, together with analytic terms that control the reference state and growth. The learned contribution combines a common family of paired directional features with one of two monotone convex cores: an input-convex neural network (ICNN) or an input-convex Kolmogorov--Arnold network (ICKAN). These cores provide alternative representations while preserving the same feature construction and mechanical guarantees.

The remainder of this section makes these distinctions precise. We first state the mechanical requirements and separate potential structure from finite-strain stability. We then define the reference model, construct the paired features and convex cores, and establish the properties of the complete energy. Throughout, $\F\in\GL$, $J=\det\F>0$, and derivatives with respect to deformation are taken with the learned parameters fixed.
\subsection{Admissibility criteria and scope of the guarantees}
\label{sec:requirements}
A small prediction error is not sufficient to define an admissible hyperelastic surrogate. We organize the construction around six criteria, following the distinctions made in constitutive PANN formulations~\cite{klein2022,linden2023}. The enumeration is an organizing device for this paper, not a universal checklist or a claim that polyconvexity is necessary for every physical material:
\begin{enumerate}
 \item[(C1)] \emph{Potential structure:} stress is the derivative of a single-valued stored energy in work-conjugate variables.
 \item[(C2)] \emph{Objectivity:} a superposed rigid rotation does not change that energy.
 \item[(C3)] \emph{Stress symmetry and anisotropy:} the material stress is symmetric, while directional dependence is allowed without imposing an unverified material symmetry group.
 \item[(C4)] \emph{Reference normalization:} the undeformed state has zero energy and zero stress.
 \item[(C5)] \emph{Polyconvexity:} the energy admits a convex representation in deformation minors, providing a sufficient rank-one convexity condition.
 \item[(C6)] \emph{Growth:} energy diverges under volume collapse, unbounded volume expansion, and unbounded distortion at bounded positive volume.
\end{enumerate}
The ICNN and ICKAN constructions below satisfy these criteria for every finite parameter set obeying their architectural constraints. Global nonnegative energy is an additional condition checked separately for the selected parameters. Constitutive identities, sufficient bounds, and sampled numerical checks consequently play different roles.

\paragraph{Potential structure (C1).}
For a differentiable hyperelastic energy, stress is its derivative in work-conjugate variables. In the convention of Eq.~\eqref{eq:voigt},
\begin{equation}
 \ssv(\ee)=\nabla_{\ee}W(\ee),\qquad
 \oint_{\Gamma}\ssv\cdot\dd\ee=0 .
 \label{eq:conservative_requirement}
\end{equation}
For the isothermal purely elastic setting considered here, (C1) also gives $\ssv\cdot\dot{\ee}-\dot W=0$: the mechanical dissipation vanishes. The closed-path identity is exact for a single-valued differentiable potential. On a simply connected strain domain, a continuously differentiable stress map has such a potential if its Jacobian is symmetric. Direct stress regression need not impose the cross-derivative identities
$\partial s_i/\partial e_j=\partial s_j/\partial e_i$~\cite{asad2022}. A small pointwise stress error alone does not establish them.

\paragraph{Objectivity and the three meanings of symmetry (C2)--(C3).}
For a superposed proper rotation $\bm Q$, objectivity requires
\begin{equation}
 W(\bm Q\F)=W(\F),\qquad
 \bm P(\bm Q\F)=\bm Q\bm P(\F).
\end{equation}
Using $\C=\F^T\F$ as an input to an energy enforces the first identity; differentiating the energy then gives the stress transformation. A direct map from $\E$ to the material stress $\Sg$ can also be objective without possessing a potential. Objectivity therefore does not distinguish Regression from Free in this study. The symmetry $\Sg=\Sg^T$ follows from differentiation with respect to symmetric strain for the energy models and is built into the three-component output of Regression. It implies a symmetric Cauchy stress $\bm\sigma=J^{-1}\F\Sg\F^T$, as required by angular-momentum balance in the present continuum.

Neither stress symmetry nor objectivity imposes a material symmetry group. A right action $\F\mapsto\F\bm Q_m$ leaves the energy unchanged only when $\bm Q_m$ belongs to that group. Directional features are fixed in the reference material and are unchanged by a superposed spatial rotation. Finally, the constitutive tangent $\bm D=\partial\ssv/\partial\ee\in\R^{3\times3}$ measures how the three stress components change with engineering strain, with entries $D_{ij}=\partial s_i/\partial e_j$. For an energy-based law it is the Hessian $\partial^2_{\ee\ee}W$. Its \emph{major symmetry}, $\bm D=\bm D^T$, is a third property: it follows from a twice differentiable scalar energy, not merely from a symmetric stress tensor.

\paragraph{Reference normalization (C4).}
Reference normalization imposes
\begin{equation}
 W(\I)=0,\qquad \bm P(\I)=\bm0.
\end{equation}
The first condition fixes an arbitrary additive energy constant. The second is a substantive stress condition and depends on how the energy is corrected. Subtracting a linear function of $\E$ preserves a potential but need not preserve polyconvexity in $\F$. An affine correction in independent deformation minors, in contrast, preserves convexity in those variables. This is the representation introduced next and used by our normalization.

\paragraph{Polyconvexity and finite-strain stability (C5).}
Potential structure does not impose a stability condition. The Legendre--Hadamard condition for a twice differentiable energy is
\begin{equation}
 \partial^2_{\F\F}W(\F)[\bm a\otimes\bm b,\bm a\otimes\bm b]\geq0
 \quad\hbox{for all }\bm a,\bm b .
 \label{eq:lh_background}
\end{equation}
Here $\bm a,\bm b\in\R^2$ are arbitrary vectors, $(\bm a\otimes\bm b)_{ij}=a_i b_j$ is their outer product, and $t\in\R$ parametrizes the perturbation. The condition tests the curvature of the scalar function $t\mapsto W(\F+t\bm a\otimes\bm b)$ at $t=0$: the energy must not curve downward along a rank-one perturbation. Strict positivity for nonzero directions is a stronger statement than this nonnegative condition. In two dimensions, a sufficient condition is a representation
\begin{equation}
 W(\F)=\mathcal P(\F,\det\F),\qquad
 \mathcal P\text{ convex in its independent arguments}.
 \label{eq:polyconvex_definition}
\end{equation}
Here $\mathcal P$ is defined on the convex set $\R^{2\times2}\times(0,\infty)$, treating its matrix and scalar arguments as independent; only its evaluation in $W$ sets the second argument equal to $\det\F$. This distinction is essential: we do not require $W$ to be convex in $\F$ itself. The cofactor is linear in the entries of a $2\times2$ matrix, so it need not be introduced as an additional independent argument. Along a rank-one perturbation, both $\F$ and $\det\F$ vary affinely with the perturbation parameter. The resulting path is therefore a straight line in the arguments of $\mathcal P$, whose convexity gives Eq.~\eqref{eq:lh_background}. Appendix~\ref{app:finite_strain_curvature} supplies the derivation. Testing finitely many rank-one directions cannot establish polyconvexity.

The standard three-dimensional definition uses $(\F,\cof\F,\det\F)$~\cite{ball1977,Ciarlet1988}. The guarantee established here is two-dimensional; it does not establish an admissible three-dimensional extension with arbitrary out-of-plane deformation.

Convexity in $\E$ is a different property because $\E=\tfrac12(\F^T\F-\I)$ depends quadratically on $\F$. Consequently, the energy curvature along a linear perturbation of $\F$ includes both a strain-tangent contribution and a contribution from the curvature of the strain path, weighted by the current stress. The latter can be negative in compression, so a positive-semidefinite strain Hessian $\bm D$ alone does not guarantee Eq.~\eqref{eq:lh_background}~\cite{asad2022,GaoNeffRoventaThiel2017}. Convexity in $\C$ and in $\E$ are equivalent up to their affine change of variables, but neither implies polyconvexity. Appendix~\ref{app:finite_strain_curvature} derives this stress-dependent curvature and gives an analytical counterexample.

\paragraph{Growth and the limits of the guarantees (C6).}
Finally, the inequalities $W(\F)\geq0$, the barrier $W\to\infty$ as $J\to0^+$, and the existence of minimizers are not synonyms. Ciarlet's presentation of Ball's existence argument~\cite{Ciarlet1988}, especially Sections 4.9 and 7.7, includes coercivity, admissible boundary data, and finite-energy competitors in addition to polyconvexity. Neither the existence theory nor a constitutive barrier guarantees convergence of an unconstrained Newton step. In this work, architectural statements, saved-parameter bounds, and numerical observations are therefore identified separately.

\subsection{Unconstrained energy reference}
\label{sec:tiers}
The most flexible energy model used in the comparison is a smooth neural network $\widehat W_{\theta}$ acting on
\begin{equation}
 \bm x(\ee)=\big(C_{11}-1,C_{22}-1,C_{12},J-1\big)^T.
 \label{eq:unconstrained_inputs}
\end{equation}
Here $\C=\I+2\E$ and $J=\sqrt{\det\C}$, so the four objective inputs are functions of the same three engineering-strain components; they do not introduce an additional state variable. Denote the resulting raw energy by $\widetilde W(\ee)=\widehat W_{\theta}(\bm x(\ee))$. A generic network does not satisfy the reference conditions automatically, so we define the normalized energy as
\begin{equation}
 W_{\mathrm{unc}}(\ee)=\widetilde W(\ee)-\widetilde W(\bm0)
 -\partial_{\ee}\widetilde W(\bm0)\cdot\ee.
 \label{eq:unconstrained_energy}
\end{equation}
At $\ee=\bm0$, the first two terms cancel and the linear term vanishes, giving $W_{\mathrm{unc}}(\bm0)=0$. Differentiation gives $\partial_{\ee}W_{\mathrm{unc}}(\ee)=\partial_{\ee}\widetilde W(\ee)-\partial_{\ee}\widetilde W(\bm0)$, and therefore zero reference stress. Because the correction is affine in $\ee$, it does not change the Hessian of the raw network. Equation~\eqref{eq:unconstrained_energy} thus fixes the reference energy and stress while leaving the learned curvature unrestricted.

Using components of $\C$ in a material reference basis ensures objectivity under superposed spatial rotations but does not impose isotropy: the network may depend differently on $C_{11}$, $C_{22}$, and $C_{12}$. A change of material reference axes must transform these tensor components and the represented law consistently. This model is used only as the flexible energy-based reference for interpreting the constrained ICNN and ICKAN results.

\subsection{Feature design: expressiveness and admissibility}
\label{sec:feature_order}
The first design question is how to describe deformation without erasing the effective anisotropy of the RVE. Classical invariants provide the natural starting point. For a three-dimensional right Cauchy--Green tensor $\C_3$, the familiar first, second, and third principal invariants are
\begin{equation}
 I_1^{(3)}=\tr\C_3,\qquad
 I_2^{(3)}=\tfrac12\bigl[(\tr\C_3)^2-\tr(\C_3^2)\bigr],\qquad
 I_3^{(3)}=\det\C_3.
 \label{eq:classical_invariants}
\end{equation}
They describe principal stretches without recording their orientation relative to the material. An energy depending only on these scalars is therefore isotropic. In the present two-dimensional setting, the corresponding two independent principal invariants are $I_1=\tr\C$ and $\det\C=J^2$; a third independent principal invariant is not available.

To retain directional information, one can supplement the isotropic measures with a structural direction $\bm d$. The familiar anisotropic stretch invariant is
\begin{equation}
 I_4(\bm d)=\bm d^T\C\bm d
          =\C:(\bm d\otimes\bm d)=\|\F\bm d\|^2,
 \qquad \|\bm d\|=1.
 \label{eq:directional_invariant}
\end{equation}
It measures squared stretch along a material direction and remains unchanged by a superposed spatial rotation. Such directional invariants are standard ingredients of anisotropic hyperelastic energies~\cite{SchroderNeff2003,Gasser2006}. Here $\bm d$ need not represent a physical fiber: learned directions provide objective scalar descriptors of the RVE's effective anisotropy. In two dimensions, $\bm d$ is naturally accompanied by its perpendicular $\bm d^\perp$. For this orthonormal pair,
\begin{equation}
 I_4(\bm d)+I_4(\bm d^\perp)=I_1,\qquad
 \|\cof\F\,\bm d\|^2=I_4(\bm d^\perp).
 \label{eq:directional_invariant_2d}
\end{equation}
The first identity shows that the pair resolves how $I_1$ is distributed between two perpendicular directions: their sum recovers $I_1$, while their individual values retain directional information. The second shows that a cofactor-based directional measure reads the stretch in the perpendicular direction rather than introducing an independent third in-plane descriptor. These observations motivate the paired directional building block used below.

Objective directional descriptors solve only the first part of the design problem. Because the downstream core is constrained to be nondecreasing in every input, the feature coordinates also determine which deformation states the core must order. The feature map must therefore be designed together with that monotonicity constraint. Consider a feature vector $\bm z(\F)$ and a scalar core $H$ nondecreasing in each input. For any two deformation states,
\begin{equation}
 \bm z(\F_B)\geq\bm z(\F_A)\text{ componentwise}
 \quad\Longrightarrow\quad H(\bm z(\F_B))\geq H(\bm z(\F_A)).
 \label{eq:feature_ordering}
\end{equation}
This implication holds regardless of the network's width. If the complete constitutive energy is $W=H(\bm z)+W_{\rm corr}$, where $W_{\rm corr}$ collects the analytic normalization and growth terms, the ordering constrains the neural contribution $W-W_{\rm corr}$ rather than the complete energy. The feature map and analytic terms must therefore be considered together.

This observation explains why the choice of an admissible feature construction can matter more than the capacity of its convex core. Changing the number of features does not change the number of independent strains: the features remain functions of the same three engineering-strain components. They embed that state into variables on which the network is constrained to be monotone and convex. Additional features can enlarge the admissible response family, but their number alone does not prove completeness.

We therefore seek features that retain directional stretch, reshape the componentwise ordering seen by the monotone core, and remain convex in the independent variables $(\F,J)$. The paired construction below achieves these three objectives by coupling directional stretch to volume through admissible determinant exponents. A feature can consequently decrease along a path when its denominator grows sufficiently, even if a numerator increases. Its reference derivative is also controlled analytically, which will allow the reference stress to be removed without sacrificing convexity.

\subsubsection{Paired directional features}
Let $m$ denote the number of scalar features. For each $i=1,\ldots,m$, let $\bm d_i=(\cos\theta_i,\sin\theta_i)^T$ and $\bm d_i^\perp=(-\sin\theta_i,\cos\theta_i)^T$. Define
\begin{equation}
 \begin{aligned}
 z_i(\F,J)
 &=\frac{I_4(\bm d_i)^{p_i}}{p_i J^{b_i}}
          +\frac{I_4(\bm d_i^\perp)^{q_i}}{q_i J^{c_i}}\\
 &=\frac{\|\F\bm d_i\|^{2p_i}}{p_i J^{b_i}}
          +\frac{\|\F\bm d_i^\perp\|^{2q_i}}{q_i J^{c_i}},
 \end{aligned}
 \label{eq:features}
\end{equation}
where $\theta_i$ sets the orientation of the reference directions, $p_i,q_i$ are the powers of the squared-stretch invariants, and $b_i,c_i$ control each contribution's coupling to the area ratio $J$. These roles provide expressiveness, but the exponents cannot be chosen arbitrarily if the feature is to remain convex. They satisfy
\begin{equation}
 p_i,q_i\geq\tfrac12,\qquad
 0\leq b_i\leq2p_i-1,\qquad 0\leq c_i\leq2q_i-1.
 \label{eq:featureconstraints}
\end{equation}
The lower bounds on $p_i$ and $q_i$ make the corresponding powers of the directional norms convex, while the upper bounds on $b_i$ and $c_i$ prevent their coupling with $J$ from introducing a negative-curvature direction. Appendix~\ref{app:convex} derives these conditions from the Hessian of a single stretch--determinant contribution.
In the learned-feature variants, $\theta_i$ and the four exponent parameters are optimized, with $p_i,q_i,b_i,c_i$ subject to Eq.~\eqref{eq:featureconstraints}; the fixed-feature controls of Section~\ref{sec:constitutive_training} hold the same initialized feature parameters constant while training the core. We use $\operatorname{softplus}(x)=\log(1+e^x)$ and $\operatorname{sigmoid}(x)=(1+e^{-x})^{-1}$, whose ranges for $x\in\R$ are $(0,\infty)$ and $(0,1)$, respectively. With unconstrained parameters $\widehat p_i,\widehat q_i,\widehat b_i,\widehat c_i\in\R$, set
\begin{equation}
 \begin{aligned}
 p_i&=\tfrac12+\operatorname{softplus}(\widehat p_i),
 &b_i&=(2p_i-1)\operatorname{sigmoid}(\widehat b_i),\\
 q_i&=\tfrac12+\operatorname{softplus}(\widehat q_i),
 &c_i&=(2q_i-1)\operatorname{sigmoid}(\widehat c_i).
 \end{aligned}
 \label{eq:feature_parameterization}
\end{equation}
The angles require no admissibility bound: for any fixed $\theta_i$, the maps $\F\mapsto\F\bm d_i$ and $\F\mapsto\F\bm d_i^\perp$ remain linear in $\F$. The transformations in Eq.~\eqref{eq:feature_parameterization} preserve the exponent inequalities during training, so a penalty loss is not needed to enforce them. For finite real values of the unconstrained parameters, the bounds in Eq.~\eqref{eq:featureconstraints} hold strictly; boundary values are approached as limits. Each $z_i$ is one scalar feature containing two perpendicular contributions. The two contributions need not have equal exponents. Consequently, pairing perpendicular directions does not impose invariance under their interchange.

The construction was also designed to simplify reference normalization. Its paired and scaled contributions make the reference strain gradient of every feature isotropic. Consequently, the neural contribution has an isotropic reference stress, which a single scalar determinant correction can cancel in Section~\ref{sec:energy_admissibility}. The following result gathers the objectivity, convexity, and reference-state properties needed there.

\begin{proposition}
\label{prop:paired_features}
Under the constraints in Eq.~\eqref{eq:featureconstraints}, each feature in Eq.~\eqref{eq:features} is objective and convex in the independent variables $(\F,J)$ on $J>0$. Its value and engineering-strain gradient at the identity, evaluated with $J=\det\F$, are
\begin{equation}
 z_i^0=\frac1{p_i}+\frac1{q_i},\qquad
 \partial_{\ee}z_i\big|_{\bm0}=\rho_i(1,1,0)^T,
 \qquad \rho_i=2-\frac{b_i}{p_i}-\frac{c_i}{q_i}.
 \label{eq:referencefeatures}
\end{equation}
\end{proposition}
The proof is given in Appendix~\ref{app:convex}. The mechanism behind the reference gradient is that $1/p_i$ and $1/q_i$ cancel the stretch-power multipliers upon differentiation. The two perpendicular directional contributions then sum to $2\I$, while the determinant terms contribute $-(b_i/p_i+c_i/q_i)\I$. Thus $\rho_i$ is the scalar amplitude of an isotropic reference gradient: its two normal components are equal and its shear component is zero. This reference property does not imply isotropy of the response away from the reference state.

The raw features are now mechanically admissible, but their numerical magnitudes may differ. Before passing them to the core, we therefore use $y_i=(z_i-z_i^{\rm cen})/s_i$, where $z_i^{\rm cen}$ is a centering value and $s_i>0$ is a scaling factor. Both are independent of deformation for any fixed parameter set. This positive affine conditioning preserves both convexity in independent $(\F,J)$ and the componentwise ordering required by the monotone core. The choices used in the constitutive assessments are specified in Section~\ref{sec:constitutive_training}.

This feature design has a precise but limited purpose: it supplies objective, expressive, and admissible inputs to the constrained core. Learning the directions does not amount to identifying or enforcing a complete material symmetry group, and no completeness or universal-approximation theorem is asserted for a finite collection of these features followed by a monotone convex core.

\subsection{Monotone ICNN and integrated-hat ICKAN cores}
\label{sec:convex_cores}
The feature construction of Section~\ref{sec:feature_order} resolves the input side of the constitutive design: each $z_i$ is objective and convex in the independent variables $(\F,J)$. The remaining question is how a learned core may combine these inputs without undoing that convexity. The core receives the conditioned vector $\bm y=(y_1,\ldots,y_m)^T$, with $y_i=(z_i-z_i^{\rm cen})/s_i$. Because $s_i>0$, this affine conditioning preserves the convexity and componentwise ordering of the physical features.

A convex core alone is not sufficient. Composition with the convex feature map preserves convexity only when the core is also componentwise nondecreasing~\cite{klein2022,linden2023}. We enforce these same two properties in two different ways. The ICNN learns nonnegative linear maps and applies fixed softplus activations at its hidden nodes. The ICKAN instead learns a convex, nondecreasing scalar function on every edge and sums the incoming contributions at each node. The following definitions make these two mechanisms precise; Figure~\ref{fig:convex_core_architectures} then compares their multilayer architectures.

\begin{figure}[tbp]
 \centering% General multilayer schematic of the deployed monotone core families.
% Ellipses denote repeated hidden layers; widths and curves are illustrative.
% Purple paths bundle nonnegative linear maps of the entire original input.
\begin{tikzpicture}[x=1cm,y=1cm,font=\small,
  weight/.style={draw=blue!70!black,line width=.65pt},
  connection/.style={draw=black!60,line width=.6pt},
  inputmap/.style={draw=violet!65!black,dashed,line width=.65pt},
  maparrow/.style={inputmap,-{Latex[length=1.4mm]}},
  dot/.style={circle,fill=black,inner sep=0pt,minimum size=4pt},
  callout/.style={-{Latex[length=1.8mm]},draw=black,line width=.75pt},
  annotation/.style={font=\small\itshape,align=center,text width=3.15cm,inner sep=2pt},
  pics/softplus/.style={code={
    \draw[draw=red!65!black,fill=white,line width=.45pt]
      (-.44,-.32) rectangle (.44,.32);
    \draw[black!25,line width=.2pt]
      (-.33,-.22) -- (.34,-.22) (-.33,-.22) -- (-.33,.24);
    \draw[red!70!black,line width=.7pt]
      plot[domain=-3:3,samples=35]
      ({.11*\x},{-.22+.42*(ln(1+exp(\x))-ln(1+exp(-3)))/3});
  }},
  pics/hatcurve/.style args={#1/#2}{code={
    \draw[draw=teal!70!black,fill=white,line width=.4pt]
      (-.44,-.32) rectangle (.44,.32);
    \draw[black!25,line width=.2pt]
      (-.33,-.22) -- (.34,-.22) (-.33,-.22) -- (-.33,.24);
    \draw[teal!75!black,line width=.6pt]
      plot[domain=-1.5:1.5,samples=41]
      ({.22*\x},{-.22+.42*(#1*(\x+1.5)
        +#2*(max(\x+1,0)^3-2*max(\x,0)^3+max(\x-1,0)^3)/6)
        /(3*#1+1.5*#2)});
  }}
]
 \path[use as bounding box] (0,0) rectangle (16.2,9.45);
 \fill[blue!3] (1.6,8.75) rectangle (8.9,9.45);
 \fill[teal!4] (8.9,8.75) rectangle (16.2,9.45);
 \fill[black!2] (0,1.1) rectangle (1.6,9.45);
 \draw[black!65,line width=.5pt] (0,0) rectangle (16.2,9.45);
 \foreach \tablex in {1.6,8.9}
   \draw[black!65,line width=.5pt] (\tablex,1.1) -- (\tablex,9.45);
 \foreach \tabley in {1.1,7.2,8.75}
   \draw[black!65,line width=.5pt] (0,\tabley) -- (16.2,\tabley);
 \node[font=\small\bfseries] at (.8,9.1) {Core};
 \node[font=\small\bfseries] at (5.25,9.1) {Monotone ICNN};
 \node[font=\small\bfseries] at (12.55,9.1) {Monotone integrated-hat ICKAN};
 \node[align=center,font=\footnotesize] at (.8,7.98) {Layer\\operation};
 \node[align=center,font=\footnotesize] at (.8,4.15) {Multilayer\\core};
 \node[align=center,text width=6.6cm] at (5.25,7.98)
   {\textcolor{blue!70!black}{Learned nonnegative linear weights} + biases\\[3pt]
    $\longrightarrow$ \textcolor{red!70!black}{fixed softplus at hidden nodes}};
 \node[align=center,text width=6.6cm] at (12.55,7.98)
   {\textcolor{teal!75!black}{One learned scalar function per edge}\\[2pt]
    \textcolor{teal!75!black}{(convex and nondecreasing)}\\[2pt]
    $\longrightarrow$ sum at nodes + biases};

 % ICNN: information flows from the conditioned input at the top to the
 % scalar output at the bottom, consistently with the paper overview.
 \begin{scope}[xshift=1.6cm,yshift=1.35cm]
  \node[anchor=west,font=\small\bfseries] at (.12,5.55) {(a)};
  \node[font=\footnotesize,text=violet!65!black] (icinputvector) at (1.6,5.62)
    {input $\bm y$};
  \foreach \level in {4.05,1.70}
    \filldraw[red!65!black,fill=red!3,line width=.75pt]
      (.02,\level-.25) rectangle (3.18,\level+.25);
  \node[dot] (icinput1) at (.9,5.10) {};
  \node[dot] (icinputm) at (2.3,5.10) {};
  \node[font=\footnotesize] at (.9,5.34) {$y_1$};
  \node[font=\footnotesize] at (2.3,5.34) {$y_m$};
  \node[font=\footnotesize] at (1.6,5.10) {$\cdots$};
  \node[dot] (icoutput) at (1.6,.55) {};
  \node[font=\footnotesize] at (1.6,.12) {$\widehat H(\bm y)$};
  \foreach \j in {0,...,2} {
    \pgfmathsetmacro{\hiddenx}{.35+1.25*\j}
    \draw[weight] (icinput1) -- (\hiddenx,4.05);
    \draw[weight] (icinputm) -- (\hiddenx,4.05);
    \draw[weight] (\hiddenx,1.70) -- (icoutput);
    \foreach \k in {0,...,2} {
      \pgfmathsetmacro{\nextx}{.35+1.25*\k}
      \draw[weight] (\hiddenx,4.05) -- (\nextx,1.70);
    }
  }
  \fill[white] (.02,2.70) rectangle (3.23,3.13);
  \foreach \hiddenx in {.35,1.6,2.85}
    \node at (\hiddenx,2.91) {$\vdots$};

  % Direct input map to the output, and bundled maps into later preactivations.
  \draw[inputmap,rounded corners=2pt] (icinputvector.east)
    -- (3.7,5.62) -- (3.7,.55);
  \draw[maparrow] (3.7,.55) -- (icoutput.east);
  \draw[inputmap] (3.7,2.22) -- (.35,2.22);
  \foreach \hiddenx in {.35,1.6,2.85}
    \draw[maparrow] (\hiddenx,2.22) -- (\hiddenx,1.95);
  \node[fill=white,inner sep=1pt,text=violet!65!black] at (3.7,2.91) {$\vdots$};
  \foreach \j in {0,...,2} {
    \pgfmathsetmacro{\hiddenx}{.35+1.25*\j}
    \pic[scale=.58] at (\hiddenx,4.05) {softplus};
    \pic[scale=.58] at (\hiddenx,1.70) {softplus};
  }
  \node[annotation,text=red!70!black] (icfixed) at (5.65,4.15)
    {Nonlinear, fixed\\softplus activations};
  \node[annotation,text=blue!70!black] (icweights) at (5.65,3.05)
    {Learned nonnegative\\linear weights};
  \node[annotation,text=violet!65!black] (icmaps) at (5.65,1.25)
    {Original input to later\\hidden layers and output};
  \draw[callout] (3.21,4.05) to[out=0,in=195] (icfixed.west);
  \draw[callout] (3.16,3.18) to[out=0,in=180] (icweights.west);
  \draw[callout] (3.77,1.25) -- (icmaps.west);
 \end{scope}

 % ICKAN: the same top-to-bottom convention; every displayed edge contains
 % its own learned scalar function before the receiving sum node.
 \begin{scope}[xshift=8.9cm,yshift=1.35cm]
  \node[anchor=west,font=\small\bfseries] at (.12,5.55) {(b)};
  \node[font=\footnotesize,text=violet!65!black] (kaninputvector) at (1.6,5.62)
    {input $\bm y$};
  \foreach \level in {4.55,2.25,1.05}
    \filldraw[teal!70!black,fill=teal!3,line width=.75pt]
      (0,\level-.22) rectangle (3.26,\level+.22);
  \node[dot] (kaninput1) at (.9,5.10) {};
  \node[dot] (kaninputm) at (2.3,5.10) {};
  \node[font=\footnotesize] at (.9,5.34) {$y_1$};
  \node[font=\footnotesize] at (2.3,5.34) {$y_m$};
  \node[font=\footnotesize] at (1.6,5.10) {$\cdots$};
  \node[dot] (kanoutput) at (1.6,.55) {};
  \node[font=\footnotesize] at (1.6,.12) {$\widehat H(\bm y)$};
  \foreach \j in {0,...,2} {
    \pgfmathsetmacro{\hiddenx}{.35+1.25*\j}
    \node[dot] (kanfirst\j) at (\hiddenx,3.95) {};
    \node[dot] (kanlast\j) at (\hiddenx,1.70) {};
  }
  \foreach \j in {0,...,2} {
    \pgfmathsetmacro{\hiddenx}{.35+1.25*\j}
    \draw[connection] (kanlast\j) -- (\hiddenx,1.05) -- (kanoutput);
    \foreach \r in {0,1} {
      \pgfmathsetmacro{\edgex}{.27+.53*(2*\j+\r)}
      \ifnum\r=0
        \draw[connection] (kaninput1) -- (\edgex,4.55) -- (kanfirst\j);
      \else
        \draw[connection] (kaninputm) -- (\edgex,4.55) -- (kanfirst\j);
      \fi
    }
    \foreach \k in {0,...,2} {
      \pgfmathsetmacro{\edgex}{.16+.36*(3*\j+\k)}
      \draw[connection] (kanfirst\k) -- (\edgex,2.25) -- (kanlast\j);
    }
  }
  \fill[white] (.02,2.78) rectangle (3.27,3.20);
  \foreach \hiddenx in {.35,1.6,2.85}
    \node at (\hiddenx,2.99) {$\vdots$};
  \foreach \j in {0,...,2} {
    \pgfmathsetmacro{\hiddenx}{.35+1.25*\j}
    \pgfmathsetmacro{\slope}{.08+.25*\j}
    \pgfmathsetmacro{\curvature}{1.5-.45*\j}
    \pic[scale=.45] at (\hiddenx,1.05) {hatcurve=\slope/\curvature};
    \foreach \r in {0,1} {
      \pgfmathsetmacro{\edgex}{.27+.53*(2*\j+\r)}
      \pgfmathsetmacro{\slope}{.08+.15*\j+.2*\r}
      \pgfmathsetmacro{\curvature}{1.4-.3*\j-.2*\r}
      \pic[scale=.39] at (\edgex,4.55) {hatcurve=\slope/\curvature};
    }
    \foreach \k in {0,...,2} {
      \pgfmathsetmacro{\edgex}{.16+.36*(3*\j+\k)}
      \pgfmathsetmacro{\slope}{.08+.11*\j+.14*\k}
      \pgfmathsetmacro{\curvature}{1.5-.23*\j-.25*\k}
      \pic[scale=.30] at (\edgex,2.25) {hatcurve=\slope/\curvature};
    }
  }
  \draw[inputmap,rounded corners=2pt] (kaninputvector.east)
    -- (3.7,5.62) -- (3.7,.55);
  \draw[maparrow] (3.7,.55) -- (kanoutput.east);
  \node[annotation,text=teal!75!black] (kanfunctions) at (5.65,4.42)
    {One learned convex,\\nondecreasing function\\per edge};
  \node[annotation] (kansums) at (5.65,3.13) {Sum operation\\at nodes};
  \node[annotation,text=violet!65!black] (kanmap) at (5.65,1.25)
    {Linear map from\\input to output};
  \draw[callout] (3.29,4.55) to[out=0,in=170] (kanfunctions.west);
  \draw[callout] (3.0,3.95) to[out=0,in=130] (kansums.west);
  \draw[callout] (3.77,1.25) -- (kanmap.west);
 \end{scope}

 % Colour legend and the sign constraints of these monotone implementations.
 \draw[weight,line width=1pt] (.35,.78) -- (.85,.78);
 \node[anchor=west,font=\footnotesize] at (.98,.78) {Learned linear weights};
 \pic[scale=.40] at (5.05,.78) {softplus};
 \node[anchor=west,font=\footnotesize] at (5.42,.78) {Fixed softplus};
 \pic[scale=.40] at (9.05,.78) {hatcurve=.1/1.3};
 \node[anchor=west,font=\footnotesize] at (9.42,.78) {Learned spline functions};
 \draw[inputmap,line width=1pt] (13.35,.78) -- (13.85,.78);
 \node[anchor=west,font=\footnotesize] at (13.98,.78) {Input maps};
 \node[font=\footnotesize] at (8.1,.29)
   {Glyph counts show operator placement, not model size; all linear weights are nonnegative.};
\end{tikzpicture}

 \caption{Multilayer monotone cores mapping the conditioned paired features $\bm y$ to the scalar contribution $\widehat H(\bm y)$. (a) The ICNN learns nonnegative linear maps and applies fixed softplus activations at hidden nodes. (b) The ICKAN learns one convex, nondecreasing scalar function per edge and sums them at nodes. Purple dashed paths denote nonnegative maps of the complete input; ellipses denote intervening layers. Glyph counts indicate where each operation acts, not deployed widths or parameter counts. Unrestricted biases are omitted, and the diagram represents the neural cores rather than the complete energy.}
 \label{fig:convex_core_architectures}
\end{figure}

\subsubsection{Monotone ICNN}
\label{sec:icnn_core}
Let $L$ denote the number of hidden layers and $\bm a^{(\ell)}$ the vector of activations in layer $\ell$, for $\ell=1,\ldots,L$. We write $\widehat H(\bm y)$ for the scalar neural core acting on the conditioned features $\bm y$; the complete energy additionally includes the reference-normalization and growth terms of Section~\ref{sec:energy_admissibility}. The deployed ICNN has the form
\begin{align}
 \bm a^{(1)}&=\phi(\bm U_1\bm y+\bm b_1),\nonumber\\
 \bm a^{(\ell)}&=\phi(\bm W_\ell\bm a^{(\ell-1)}+\bm U_\ell\bm y+\bm b_\ell),
 \qquad \ell=2,\ldots,L,\nonumber\\
 \widehat H(\bm y)&=\bm v^T\bm a^{(L)}+\bm w^T\bm y+b_{\rm out},
 \label{eq:monotone_icnn}
\end{align}
where $\phi=\operatorname{softplus}$ is applied componentwise. Its derivatives satisfy $\phi'(x)=\operatorname{sigmoid}(x)>0$ and $\phi''(x)=\operatorname{sigmoid}(x)[1-\operatorname{sigmoid}(x)]>0$. Thus the activation is smooth, increasing, and convex; its smoothness also makes the core $C^2$, as required for a continuous constitutive tangent. The matrices $\bm U_\ell$ and $\bm W_\ell$ contain the weights from $\bm y$ and the preceding hidden layer, respectively, and $\bm b_\ell$ is the layer bias vector. The output weight vectors $\bm v$ and $\bm w$ act on $\bm a^{(L)}$ and $\bm y$, respectively; $b_{\rm out}$ is an additive scalar output bias.

Every entry of $\bm W_\ell,\bm U_\ell,\bm v,\bm w$ is nonnegative, enforced by softplus transformations of unconstrained weight parameters. These signs play two related but distinct roles. The matrices $\bm W_\ell$ and the output vector $\bm v$ multiply convex hidden activations, so their nonnegativity preserves convexity as well as monotonicity. The maps $\bm U_\ell\bm y$ and $\bm w^T\bm y$ are affine and therefore do not affect convexity for either sign, but nonnegative $\bm U_\ell$ and $\bm w$ are required for componentwise monotonicity. Biases are unrestricted because adding constants to a preactivation preserves both properties. Hidden biases can change the core's derivatives with respect to $\bm y$, whereas $b_{\rm out}$ only shifts $\widehat H$ and leaves its derivatives unchanged.

The preservation argument proceeds layer by layer. The first hidden layer is convex because softplus is convex and its argument is affine, and it is nondecreasing because $\bm U_1$ is nonnegative. Assuming that $\bm a^{(\ell-1)}$ is convex and nondecreasing, the nonnegative map $\bm W_\ell\bm a^{(\ell-1)}$, the nonnegative affine input map, and composition with $\phi$ preserve both properties in layer $\ell$. The same argument applies to the output, proving by induction that $\widehat H$ is convex and componentwise nondecreasing. This is the ICNN composition argument of Amos et al.~\cite{amos2017}, supplemented here by nonnegative direct-input maps to guarantee monotonicity.

\subsubsection{Monotone integrated-hat ICKAN}
\label{sec:ickan_core}
Unlike the ICNN, which applies a fixed activation after an affine weighted sum at each hidden node, the ICKAN assigns a learned scalar function to each connection (edge) between layers. Each node sums the transformed contributions from its incoming connections. To make the relation between these edge functions and the scalar core explicit, set $\bm h^{(0)}=\bm y$ and let $n_\ell$ be the width of layer $\ell$, with $n_0=m$. The hidden layers and scalar output are
\begin{align}
 h_j^{(\ell)}
 &=\sum_{i=1}^{n_{\ell-1}}f_{ji}^{(\ell)}\!\left(h_i^{(\ell-1)}\right)
   +b_j^{(\ell)},
 &&j=1,\ldots,n_\ell,\quad \ell=1,\ldots,L,\nonumber\\
 \widehat H(\bm y)
 &=\sum_{i=1}^{n_L}f_i^{\rm out}\!\left(h_i^{(L)}\right)
   +\bm w^T\bm y+b_{\rm out},
 &&\bm w\geq\bm0.
 \label{eq:monotone_ickan}
\end{align}
Thus $f_{ji}^{(\ell)}$ transforms one scalar carried by an edge, $h_j^{(\ell)}$ is the sum formed at a node, and $\widehat H(\bm y)$ is the complete ICKAN core. The nonnegative linear skip $\bm w^T\bm y$ connects the conditioned input directly to the scalar output.
\begin{samepage}
It remains to specify the functions appearing on the edges. For any $f_{ji}^{(\ell)}$ or $f_i^{\rm out}$ in Eq.~\eqref{eq:monotone_ickan}, let $x$ denote its scalar input. Each edge uses the common form
\begin{equation}
 f(x)=a x+\sum_k c_k B_k(x),\qquad a,c_k\geq0,
 \label{eq:splineedge}
\end{equation}
where $a$ is the learned linear slope, $c_k$ are learned weights of the spline bases $B_k$, and $k$ indexes these bases. Softplus transformations of unconstrained parameters impose the signs of $a$ and $c_k$. Additive constants are represented once, through the unrestricted node biases in Eq.~\eqref{eq:monotone_ickan}, rather than redundantly on every incoming edge.
\end{samepage}

The remaining design task is to learn the shape of each edge without allowing negative slope or curvature. Rather than check these inequalities at sampled inputs after training, we parameterize the edge curvature directly. Each $B_k''$ is a nonnegative triangular hat centered at a knot $\kappa_k$. The knots form a fixed uniform grid with spacing $h>0$; neither their locations nor the spacing are optimized. Writing $s=(x-\kappa_k)/h$, one curvature basis is
\begin{equation}
 B_k''(x)=\max(1-|s|,0).
 \label{eq:hat_curvature}
\end{equation}
Integrating twice from the left conditions $B_k(x)=B_k'(x)=0$ for $x\leq\kappa_k-h$ gives
\begin{equation}
 B_k(x)=h^2\begin{cases}
 0,&s\leq-1,\\
 (s+1)^3/6,&-1<s\leq0,\\
 (-s^3+3s^2+3s+1)/6,&0<s<1,\\
 s,&s\geq1.
 \end{cases}
 \label{eq:hat}
\end{equation}
Figure~\ref{fig:integrated_hat_construction} shows both levels of the construction: integration of one curvature basis and assembly of a learned edge from several bases.
\begin{figure}[tbp]
 \centering% Analytic schematic of one integrated-hat basis and a learned edge assembled
% from several nonnegative curvature contributions.  Curves in panel (a) are
% exact for kappa=0 and h=1; panel (b) is a qualitative multi-basis schematic.
\begin{tikzpicture}[x=1cm,y=1cm,font=\small,
  axis/.style={draw=black!45,line width=.45pt},
  guide/.style={draw=black!30,dashed,line width=.4pt},
  curvature/.style={draw=orange!85!black,line width=1.25pt},
  slope/.style={draw=teal!75!black,line width=1.25pt},
  basis/.style={draw=blue!65!black,line width=1.25pt},
  total/.style={draw=black,line width=1.35pt},
  integrate/.style={-{Latex[length=1.5mm]},draw=black!60,line width=.55pt}
]
 \path[use as bounding box] (0,0) rectangle (16.2,6.35);
 \fill[blue!2] (0,0) rectangle (7.95,6.35);
 \fill[teal!2] (8.25,0) rectangle (16.2,6.35);
 \draw[black!55,line width=.5pt] (0,0) rectangle (7.95,6.35);
 \draw[black!55,line width=.5pt] (8.25,0) rectangle (16.2,6.35);

 \node[font=\small\bfseries] at (3.98,5.96) {(a) One integrated-hat basis};
 \node[font=\small\bfseries] at (12.23,5.96) {(b) One learned ICKAN edge};

 % Panel (a): exact normalized single-basis construction.  The larger left
 % inset keeps all row labels inside the panel.
 \begin{scope}[xshift=3.00cm,yshift=.62cm,x=1.40cm,y=.80cm]
  \foreach \yy/\lab in {4.65/{$B_k''$},2.70/{$B_k'$},.55/{$B_k$}} {
    \draw[axis] (-1.45,\yy) -- (1.55,\yy);
    \node[anchor=east] at (-1.60,\yy) {\lab};
  }
  \foreach \xx in {-1,0,1}
    \draw[guide] (\xx,.42) -- (\xx,5.55);

  \draw[curvature] (-1.45,4.65) -- (-1,4.65) -- (0,5.55)
    -- (1,4.65) -- (1.55,4.65);

  \draw[slope] (-1.45,2.70) -- (-1,2.70);
  \draw[slope] plot[domain=-1:0,samples=35]
    (\x,{2.70+.72*(\x+1)^2/2});
  \draw[slope] plot[domain=0:1,samples=35]
    (\x,{2.70+.72*(.5+\x-.5*\x^2)});
  \draw[slope] (1,3.42) -- (1.55,3.42);

  \draw[basis] (-1.45,.55) -- (-1,.55);
  \draw[basis] plot[domain=-1:0,samples=35]
    (\x,{.55+.72*(\x+1)^3/6});
  \draw[basis] plot[domain=0:1,samples=35]
    (\x,{.55+.72*(-\x^3+3*\x^2+3*\x+1)/6});
  \draw[basis] plot[domain=1:1.55,samples=15]
    (\x,{.55+.72*\x});

  \draw[integrate] (1.78,4.15) -- node[right,font=\footnotesize] {$\int\!\dd x$} (1.78,3.35);
  \draw[integrate] (1.78,2.18) -- node[right,font=\footnotesize] {$\int\!\dd x$} (1.78,1.38);
  \node[font=\footnotesize,anchor=north] at (-1,.34) {$\kappa_k-h$};
  \node[font=\footnotesize,anchor=north] at (0,.34) {$\kappa_k$};
  \node[font=\footnotesize,anchor=north] at (1,.34) {$\kappa_k+h$};
  \node[font=\footnotesize,anchor=west] at (1.58,.55) {$x$};
  \node[font=\scriptsize,text=blue!65!black,align=center] at (2.52,.82)
    {zero left tail\\linear right tail};
 \end{scope}

 % Panel (b): nonnegative learned combination of several fixed hats.
 \begin{scope}[xshift=11.20cm,yshift=.62cm,x=1.05cm,y=.80cm]
  \foreach \yy/\lab in {4.65/{$f''$},2.70/{$f'$},.55/{$f$}} {
    \draw[axis] (-2.05,\yy) -- (2.05,\yy);
    \node[anchor=east] at (-2.20,\yy) {\lab};
  }

  % Individual learned curvature contributions.
  \draw[blue!65!black,line width=.75pt]
    (-1.75,4.65) -- (-1.05,4.97) -- (-.35,4.65);
  \draw[orange!85!black,line width=.75pt]
    (-.80,4.65) -- (0,5.55) -- (.80,4.65);
  \draw[teal!75!black,line width=.75pt]
    (.15,4.65) -- (.90,5.15) -- (1.65,4.65);
  \draw[total] plot[domain=-1.85:1.75,samples=90]
    (\x,{4.65+.88*(.35*max(1-abs((\x+1.05)/.7),0)
      +max(1-abs(\x/.8),0)+.55*max(1-abs((\x-.9)/.75),0))});
  \node[font=\scriptsize,text=blue!65!black] at (-1.05,5.15) {$c_1$};
  \node[font=\scriptsize,text=orange!85!black] at (0,5.74) {$c_2$};
  \node[font=\scriptsize,text=teal!75!black] at (.90,5.35) {$c_3$};

  % Monotone accumulated slope and convex connection (schematic).
  \draw[slope] plot[smooth,tension=.65] coordinates {
    (-2.05,2.82) (-1.55,2.83) (-1.10,2.96) (-.65,3.08)
    (-.25,3.28) (.15,3.66) (.55,3.88) (.95,4.00)
    (1.35,4.16) (1.75,4.20) (2.05,4.20)};
  \draw[basis] plot[smooth,tension=.55] coordinates {
    (-2.05,.55) (-1.55,.57) (-1.05,.62) (-.55,.72)
    (-.05,.90) (.45,1.16) (.95,1.50) (1.45,1.85)
    (2.05,2.15)};

  \draw[integrate] (2.25,4.15) -- node[right,font=\footnotesize] {$\int\!\dd x$} (2.25,3.35);
  \draw[integrate] (2.25,2.18) -- node[right,font=\footnotesize] {$\int\!\dd x$} (2.25,1.38);
  \node[font=\footnotesize,anchor=west] at (2.08,.55) {$x$};
  \node[font=\scriptsize,align=center] at (0,-.03)
    {$a\geq0$: slope,\quad $c_k\geq0$: curvature};
 \end{scope}
\end{tikzpicture}

 \caption{Integrated-hat parameterization of an ICKAN edge. (a) A localized nonnegative curvature hat integrates to a nonnegative slope and a convex $C^2$ basis with a zero left tail and a linear right tail. (b) Nonnegative learned coefficients combine several fixed bases into a nonnegative edge curvature; integration gives a nondecreasing slope and a convex, nondecreasing connection. The curves illustrate the analytic construction; they do not represent the deployed grid or parameter count.}
 \label{fig:integrated_hat_construction}
\end{figure}

The value and first two derivatives in Eq.~\eqref{eq:hat} match at every join, and $B_k''$ vanishes at the outer joins. Thus $B_k$ is globally $C^2$. Its curvature is localized to $[\kappa_k-h,\kappa_k+h]$, but the integrated basis is not: it has a zero left tail and a linear right tail. Moreover, $B_k'(-\infty)=0$ and $B_k''\geq0$ imply $B_k'\geq0$. The complete edge therefore satisfies
\begin{equation}
 f''(x)=\sum_k c_kB_k''(x)\geq0,
 \qquad
 f'(x)=a+\sum_k c_kB_k'(x)\geq0,
 \label{eq:splineedge_guarantees}
\end{equation}
for every learned parameter set with $a,c_k\geq0$. Hence each edge is convex and nondecreasing by construction, rather than by a sampled numerical test.

Equation~\eqref{eq:monotone_ickan} sums these scalar edge functions at each node. Because every edge is convex and nondecreasing, induction across the layers shows that every $h_j^{(\ell)}$, and finally $\widehat H$, is convex and componentwise nondecreasing in $\bm y$. The nonnegative input-to-output skip preserves both properties. The fixed knot grid and finite number of bases limit where an individual edge can learn curvature: outside their combined support, the edge has linear tails. These tails are linear in the scalar input $x$, not in $\F$; the complete energy still contains nonlinear features and analytic growth terms. The grid used in the constitutive assessment is reported in Section~\ref{sec:constitutive_training}. This core is inspired by the monotonic input-convex KAN framework~\cite{thakolkaran2025}, with the integrated-hat bases and analytic tails specified here.

\subsection{The complete energy and its admissibility guarantees}
\label{sec:energy_admissibility}
Sections~\ref{sec:feature_order} and~\ref{sec:convex_cores} constructed the physical feature vector $\bm z$, its conditioned representation $\bm y(\bm z)$, and the scalar neural core $\widehat H(\bm y)$. We write
\begin{equation}
 H(\bm z)=\widehat H\bigl(\bm y(\bm z)\bigr)
 \label{eq:conditioned_neural_contribution}
\end{equation}
for the resulting learned contribution. This term represents the anisotropic response, but it does not by itself set the energy and stress references or prescribe growth at the boundary of the admissible deformation domain.

To complete the construction, we first quantify the residual stress produced by $H$ at the undeformed reference state. Let $\bm z^0=\bm z(\I,1)$. With the feature-gradient factors $\rho_i$ from Eq.~\eqref{eq:referencefeatures}, define the core sensitivities and the scalar amplitude of this residual stress as
\begin{equation}
 h_i=\partial_{z_i}H(\bm z^0)\geq0,
 \qquad
 r=\sum_i h_i\rho_i.
\end{equation}
For a fixed parameter set, $H(\bm z^0)$ and $r$ are constants with respect to deformation. The complete stored energy then combines the learned contribution with analytic reference-normalization and growth terms:
\begin{align}
 W(\F)={}&\underbrace{H(\bm z)-H(\bm z^0)}_{\text{(i) }W_{\rm learn}}
 +\underbrace{\big[-r(J-1)\big]}_{\text{(ii) }W_{\rm corr}}\nonumber\\
 &+\underbrace{\alpha(J-1-\log J)+\frac\beta2(J-1)^2}_{\text{(iii) }W_{\rm vol}}\nonumber\\
 &+\underbrace{\varepsilon\left(\frac{\tr\C}{2}-1-\log J\right)}_{\text{(iv) }W_{\rm dist}},
 \qquad \alpha,\beta,\varepsilon>0.
 \label{eq:energy}
\end{align}
Here $J=\det\F$, $\C=\F^T\F$, and $\bm z=\bm z(\F,J)$. The positive energy-density coefficients $\alpha$, $\beta$, and $\varepsilon$ weight the analytic growth terms. For either core, the composition in Eq.~\eqref{eq:conditioned_neural_contribution} is twice continuously differentiable, convex, and componentwise nondecreasing in $\bm z$.
The underbraces name the four contributions used below: $W_{\rm learn}$ carries the learned anisotropic response, while $W_{\rm corr}$, $W_{\rm vol}$, and $W_{\rm dist}$ impose the reference and growth behavior.

Each contribution solves a distinct constitutive problem:
\begin{itemize}
 \item[(i)] \textbf{Learned response and energy reference ($W_{\rm learn}$).}
 At $\F=\I$, the features equal $\bm z^0$, and hence
 \begin{equation}
  W(\I)=H(\bm z^0)-H(\bm z^0)=0.
  \label{eq:energy_reference_value}
 \end{equation}
 The remaining terms also vanish because $J=1$ and $\tr\C=2$. Since $H(\bm z^0)$ is constant with respect to deformation, its subtraction fixes the energy reference without changing the stress or tangent.

 \item[(ii)] \textbf{Reference-stress correction ($W_{\rm corr}$).}
 A constant subtraction cannot change stress. By the chain rule and Eq.~\eqref{eq:referencefeatures}, the learned contribution instead has the reference value
 \begin{equation}
  \left.\nabla_{\ee}H(\bm z)\right|_0
  =\sum_i h_i\left.\nabla_{\ee}z_i\right|_0
  =\sum_i h_i\rho_i(1,1,0)^T
  =r(1,1,0)^T.
  \label{eq:neural_reference_stress}
 \end{equation}
 The identity $J=\sqrt{(1+2E_{11})(1+2E_{22})-(2E_{12})^2}$ gives $\nabla_{\ee}J|_0=(1,1,0)^T$. Holding $r$ fixed when differentiating therefore yields
 \begin{equation}
  \left.\nabla_{\ee}\big[H(\bm z)-r(J-1)\big]\right|_0
  =r(1,1,0)^T-r(1,1,0)^T=\bm0.
  \label{eq:energy_reference_stress}
 \end{equation}
 A single scalar $r$ suffices because the paired features make only the \emph{reference} stress isotropic; the learned response away from $\F=\I$ remains anisotropic. The correction is zero in value at $J=1$ and, being affine in the independent determinant variable, preserves convexity for either sign of $r$. This affine determinant normalization follows the strategy of Linden et al.~\cite{linden2023}; Eq.~\eqref{eq:referencefeatures} makes it applicable to the present directional feature family.

 \item[(iii)] \textbf{Volumetric growth ($W_{\rm vol}$).}
 Its two parts control opposite limits: $\alpha(J-1-\log J)$ diverges as $J\to0^+$, whereas $\beta(J-1)^2/2$ grows quadratically as $J\to\infty$. At the reference,
 \begin{equation}
  W_{\rm vol}(1)=0,\qquad
  W_{\rm vol}'(J)=\alpha(1-J^{-1})+\beta(J-1),\qquad
  W_{\rm vol}'(1)=0.
 \end{equation}
 Thus this term changes neither the reference energy nor the reference stress. Both parts are needed: the quadratic term alone would remain finite at collapse, with limit $\beta/2$.

 \item[(iv)] \textbf{Distortional growth at bounded area ($W_{\rm dist}$).}
 A function of $J$ alone cannot detect a change of shape that preserves area. For example, with $\lambda>0$,
 \begin{equation}
  \F=\diag(\lambda,\lambda^{-1}),\qquad
  J=1,\qquad
  \varepsilon\left(\frac{\tr\C}{2}-1-\log J\right)
  =\frac\varepsilon2\left(\lambda^2+\lambda^{-2}-2\right).
 \end{equation}
 This contribution diverges as $\lambda\to\infty$, although the area remains constant. More generally, $\tr\C=\|\F\|_F^2$ supplies growth under unbounded deformation at bounded positive $J$. At the reference, the constant $-1$ sets its value to zero. Using $\C=\I+2\E$ and $\partial_{\E}J|_0=\I$, its stress contribution is
 \begin{equation}
  \left.\partial_{\E}\left[\varepsilon\left(\frac{\tr\C}{2}-1-\log J\right)\right]\right|_0
  =\varepsilon\left[\I-\left.J^{-1}\partial_{\E}J\right|_0\right]
  =\bm0.
 \end{equation}
\end{itemize}
Together, $W_{\rm learn}$ and $W_{\rm corr}$ set the learned response and its reference state, while $W_{\rm vol}$ and $W_{\rm dist}$ prescribe growth outside the sampled regime. Zero value and zero first derivative at the reference do not imply zero second derivative: the analytic terms still contribute to the reference stiffness because $J$ is nonlinear in $\ee$. The following proposition gathers the guarantees of their complete sum.

\begin{proposition}
\label{prop:complete_energy}
Let the features satisfy Eq.~\eqref{eq:featureconstraints}, let the core be either of the $C^2$ convex, componentwise nondecreasing constructions in Section~\ref{sec:convex_cores}, and let $\alpha,\beta,\varepsilon>0$. Then Eq.~\eqref{eq:energy} defines a $C^2$, objective, two-dimensional polyconvex energy on $\GL$. Its work-conjugate stress $\ssv=\nabla_{\ee}W$ is symmetric, $W(\I)=0$, $\ssv(\bm0)=\bm0$, and $W$ diverges under collapse, unbounded expansion, and unbounded distortion at bounded positive area. Thus the construction satisfies (C1)--(C6).
\end{proposition}
The proof is given in Appendix~\ref{app:complete_energy_proof}. In brief, the convex, componentwise nondecreasing core preserves feature convexity; the remaining terms are constant, affine, or convex in independent $(\F,J)$. The reference identities and growth limits then complete the argument. These guarantees hold for every admissible parameter set and therefore do not depend on the training samples or fitting error. Predictive accuracy and structural deployment are assessed separately in Sections~\ref{sec:constitutive_accuracy} and~\ref{sec:coupon}. Global nonnegative energy is not an architectural guarantee; it is audited for the saved checkpoints in Section~\ref{sec:mechanical_assessment} using the sufficient bound derived in Appendix~\ref{app:nonnegative}.

The three constitutive constructions carried into the numerical assessment are summarized in Table~\ref{tab:tiers}. The ICNN and ICKAN columns include both fixed- and learned-feature variants; learning the feature parameters changes the representation, not the guarantees established above.
\begin{table}[htbp]
\centering\small
\caption{Constitutive constructions carried into the numerical assessment. Criteria (C1)--(C6) are defined in Section~\ref{sec:requirements}; guarantees refer to the specific constructions in this section.}
\label{tab:tiers}
\begin{tabularx}{\linewidth}{lYYY}
\toprule
 & \shortstack{Unconstrained\\energy} & ICNN & ICKAN\\
\midrule
Role & Flexible hyperelastic reference & Guaranteed paired-feature energy & Guaranteed paired-feature energy\\
\midrule
Inputs & Objective measures from $(\C,J)$ & \shortstack{Paired features\\(fixed or learned)} & \shortstack{Paired features\\(fixed or learned)}\\
\midrule
Learned scalar map & Smooth energy network & \shortstack{Nonnegative weights;\\fixed activations} & \shortstack{Convex, nondecreasing\\edge functions}\\
\midrule
(C1)--(C4) & By construction & By construction & By construction\\
\midrule
(C5)--(C6) & Not imposed & By construction & By construction\\
\bottomrule
\end{tabularx}
\end{table}
\FloatBarrier
Section~\ref{sec:cell} evaluates whether the implementations reproduce these constructions and how accurately the resulting energies approximate the RVE response. Section~\ref{sec:reduced} first defines the intrusive microscopic references used later in the structural accuracy--cost comparison.
